\section*{September 7, 2026: Public closure signals before 2013}
\textit{Drafted by Codex from board records, contemporary reporting, and research reports.}

Public school-reorganization proposals reached neighborhoods in both comparison
groups before 2013. This supports taking anticipation seriously, but does not
establish that anticipated closure explains the 2011--12 increases in sale
prices. Earlier proposals often ended in reprieves or involved replacing staff
rather than closing a building. The relevant history is therefore a sequence of
changing risks and other interventions, not one common early treatment date.

Peabody, one of the 30 treated schools, had already faced closure in 2009.
The East Village Association's February 28, 2013 report recalled the earlier
proposal and successful opposition; its principal still described the current
outcome as uncertain. Yale, also treated, faced a proposed staff replacement
in 2009 that was withdrawn after hearings. These are different kinds of school
actions and should be recorded separately.\footnote{\href{https://www.eastvillagechicago.org/2013/02/peabody-school-again-closing-target.html}{East Village Association, February 28, 2013}; \href{https://abc7chicago.com/archive/6675416/}{ABC7, February 24, 2009}.}

Marconi, another treated school, was named in a January 19, 2010 public notice.
The proposal would have consolidated its program into Tilton while keeping
Marconi's building in school use. The February 24 board proceedings mark the
proposal as deferred. In March 2011, a different proposal made Marconi a receiving
school for Tilton students and proposed closing Beidler, a control school.
The union's May 2011 publication reports that both plans were abandoned after
community opposition.\footnote{\href{https://www.cpsboe.org/content/documents/feb24_2010proceedings.pdf}{CPS proceedings, February 24, 2010}, report 10-0224-EX5, printed pp.~36--37; \href{https://www.wbez.org/eight-forty-eight/2011/03/23/task-forces-recommendations-for-cps-school-closings}{WBEZ, March 23, 2011}; \href{https://www.ctulocal1.org/wp-content/uploads/2018/11/2011-05.pdf}{Chicago Union Teacher, May 2011}, pp.~10--11. The union publication is a participant's account.}

Five of the 49 control schools had an additional intervention during the period
of price divergence: Fuller, Herzl, Piccolo, Stagg, and Woodson South. The board
approved replacing their staff, including principals, on February 22, 2012,
effective June 30. The proceedings document written notices and public hearing
notices beginning around November 30, 2011. Matching their school IDs to the
saved descriptive packet identifies 968 of the 8,083 control-group sales over
2008--18, or 12.0 percent. They are therefore a meaningful part of the
transaction-weighted comparison.\footnote{\href{https://www.cpsboe.org/content/documents/feb2012proceedings.pdf}{CPS proceedings, February 22, 2012}, reports 12-0222-EX14, EX15, EX17, EX19, and EX21; \href{https://www.cpsboe.org/meetings/board-actions/24}{adopted board actions}. IDs: Fuller 609928; Herzl 609991; Piccolo 610106; Stagg 610339; Woodson South 610345. Their respective sale counts are 162, 44, 418, 283, and 61.}

The 2013 selection process itself distinguished recently reorganized schools.
The Consortium's account quotes CPS's policy of protecting schools that had
undergone a turnaround that school year. Recent intervention could therefore
affect both a school's housing-market trajectory and its probability of
remaining open. Being on the same February list does not remove that concern.
This is a reason to examine historical school actions within the existing
comparison, not evidence of the direction or size of their effects on prices.

\clearpage
\textbf{Timing of the mass-closure process.}

WBEZ reported on March 23, 2011 that interim CPS chief Terry Mazany had
publicly discussed closing 10--30 schools over the next two years. Broad
expectations of additional closures therefore predate the mass-closure list.
By September 11, 2012, WTTW was reporting a union researcher's account of rumors
that as many as 100 schools could close. That is evidence that the possibility
of a large closure wave was public, not that its eventual locations were known.
The Consortium dates the release of the 330-school underutilization list to
December 5, 2012 and the 129-school candidate list to February 13, 2013.
Contemporary reporting places the proposed closure announcements on March 21;
the board voted on May 22.\footnote{\href{https://news.wttw.com/2012/09/11/charter-schools-open-during-strike}{WTTW, September 11, 2012}; \href{https://files.eric.ed.gov/fulltext/ED589712.pdf}{UChicago Consortium, School Closings in Chicago (2018)}, printed pp.~14--15, for selection criteria and timeline. The figure labels the March announcement March 23; \href{https://www.chicagomag.com/city-life/march-2013/heres-the-list-of-chicago-public-schools-that-could-close/}{Chicago Magazine's March 21, 2013 reporting} establishes that names were already public on March 21.}

The December 2012 and February 2013 lists provide concrete dates for anticipation
of the mass closure wave. Those lists cannot themselves explain a rise during
2011 or most of 2012. The earlier school-specific proposals establish that some local
expectations could have changed sooner, but their withdrawals matter too:
a price increase following a reprieve could reflect a reduction in closure
risk. Annual medians and percentiles cannot reveal whether buyers responded
to a proposal, its withdrawal, another neighborhood change, or the mix of
properties offered for sale.

The useful next descriptive comparison is to locate the 2011--12 changes by
school site and property type, then align those changes with documented notices,
withdrawals, and actual school actions. A separate sensitivity plot can omit
the five 2012 turnaround controls while retaining the current full comparison
as the reference. Neither an earlier treatment date nor a revised control
sample has been adopted here.

This search establishes examples, not a complete action history for all 79
schools. It covered the 2009--12 closure and consolidation rounds, the 2012
turnarounds, the 2012--13 mass-closure chronology, and targeted searches for
northern treated schools. No public record located in this search establishes
a common 2011 announcement of the eventual 2013 closures. This does not establish
that rumors, private information, or other school-specific warnings were absent.
There is also no direct evidence here that housing buyers acted on the notices.
The reading has not yet connected the documented events to individual sales.

\small
Reproduction: source links above were reviewed September 7, 2026. The school
membership and sales totals use the fixed version at commit \texttt{e76f4dd}
and the \texttt{support} table in the descriptive packet's
\texttt{data\_review.rds}, linked to this logbook as
\texttt{../input/overview\_data\_review.rds}. No eligibility rule, treatment
date, regression, or source dataset changed. Compile with
\texttt{make -C logbook/code}.
\normalsize
